\section{Conclusions}
\label{sec:conclusion}

The optimal current-allocation problem becomes small only after its geometry
is respected. In the linear EESM, copper loss is a positive-definite
ellipsoid, torque is a signature-\((+,-,0)\) hyperbolic cylinder, and squared
stator voltage is a rank-two elliptic cylinder. Their common degree-two
homogeneity separates current direction from radial operating level.

Loss whitening gives every coordinate the same copper price. An orthogonal
rotation discovered directly from the torque expression then produces one
active-flux coordinate, one quadrature coordinate, and one torque-neutral
coordinate. Torque reduces to \(M=\kappa z_\psi z_q\). From that single
product follow the balanced unconstrained optimum, hyperbolic allocation
angle, projective loss circles, and conformal square-map picture. The
hyperbolic and projective coordinates are the same factor ratio in logarithmic
and ordinary form; the square map records its product and difference-square
geometry.

The strongest computational view maps every loss-normalized direction to
\((U^2/P_{\mathrm{Cu}},M/P_{\mathrm{Cu}})\). For three real current
coordinates this joint numerical range is convex. Rotating a support line
there yields a symmetric \(3\times3\) generalized eigenproblem and a bounded
scalar search. Minimum loss and maximum available torque become two radial
queries on the same boundary. The hard thermal limit remains explicit:
\(P_{\mathrm{req}}\leq P_{\mathrm{Cu},\max}\) means the requested torque is
physically admissible; otherwise the controller returns the maximum feasible
torque and a limit status.

The fixed-excitation PSM slice explains both what is inherited and what is
lost. Its PM flux is an affine offset, so the voltage cylinder cuts to an
ellipse whose exact resistive center follows a speed-independent conic path.
Resistance, speed, and saliency jointly rotate the ellipse; the nonsalient
circle case must be separated before evaluating an angle. The characteristic
current's position relative to the copper-loss circle decides whether the
ideal high-speed feasible set persists or terminates at finite speed.

Reading the value functions in reverse creates a machine-capability geometry.
For the voltage-free linear EESM, the time-valued coordinates
\(x_M=\Delta L/a\) and
\(y_M=L_{\mathrm{exc}}/\sqrt{aR_{\mathrm{exc}}}\) turn constant maximum
torque per copper loss into circles. Their radius is the exact spectral
descriptor \(2\lambda_{\max}(\widetilde H)/k_T\); their angle separates
reluctance and wound-field mechanisms without judging the sign convention of
saliency. Scaling by a declared base frequency makes the map dimensionless.
The limit \(R_{\mathrm{exc}}\to0\) sends the loss ellipsoid to a cylinder and
the capability coordinate to infinity. This exposes a missing physical cost;
it is not the PSM limit, which instead fixes excitation and its effective flux.

For a nonsalient PSM, low-speed torque supplies an exact minimum PM-flux
condition. Combining it with the ideal characteristic-current criterion
produces a transparent flux-inductance design window. Saliency adds the
normalized loading \(\sigma=\Delta L I_P/\psi_{\mathrm{pm}}\), while
\(\iota=I_{\mathrm{ch}}/I_P\) independently records access to the zero-flux
point. Voltage-aware EESM and PSM boundaries remain implicit level sets of the
same small support-eigenvalue solver. A torque-speed specification is their
intersection, and a robust design requires its entire temperature- and
saturation-dependent parameter path to remain inside that intersection.

Electromagnetic inversion does not yield a unique physical machine. The
forward map from geometry, winding, and materials to electromagnetic
parameters is generally many-to-one; the proper result is the feasible
preimage \(\mathcal G_{\mathrm{feas}}\). The analytic and spectral tests are
therefore inexpensive filters and sensitivity tools before nonlinear magnetic,
thermal, mechanical, demagnetization, and finite-element verification. They
also make machine-control co-design explicit: geometry shapes the quadratic
forms, and optimal current allocation evaluates their usable capability.

KKT equations, resultants, and PMSM quartics remain important independent
checks. They are not the preferred opening language because they hide
directional efficiency and branch geometry. Likewise, the complex square map
and stereographic projection teach useful invariants but make the full EESM
voltage relation no smaller.

For saturation, the next practical step is not to insert current-dependent
inductances into global linear formulas. It is to combine the analytic inner
direction solver with local affine flux models that include their tangent
offsets, damping, region checks, and original-model residuals. Open questions
include rigorous contraction regions for that outer loop, certified fixed-point
implementations of the support eigensolver, certified surrogate margins for
the physical-to-electromagnetic map, nonsmooth design sensitivities at
active-set changes, and conditions under which special machine symmetries
reduce the generic sextic projective contact polynomial. The same mathematical
structure has thus been read in both directions: machine to feasible currents
to optimal operation, and required operation to admissible electromagnetic
parameters to candidate physical designs.

Dynamic operation adds a third, genuinely different geometry. A reciprocal
EESM flux matrix must contain the field self inductance and the same mutual
inductance in stator and rotor flux equations. Its inverse differential
inductance maps the converter-voltage body into a local set of current
velocities; resistance, back-EMF, and cross coupling translate that body as an
electrical drift. Dynamic whitening makes ellipsoidal voltage authority
isotropic and exposes fast, medium, and slow current directions. These axes
generally differ from both loss-whitened and torque-producing axes.

No distance was introduced before its objective. Pure minimum time with an
ellipsoidal controlled-speed body and sufficiently weak drift is exactly a
Zermelo navigation problem and induces a nonreversible Randers metric. A real
inverter hexagon crossed with a separate field-voltage interval instead gives
a polyhedral Finsler gauge. Minimum control energy gives a controllability
Gramian; minimum state-loss energy gives an LQR-type boundary-value problem;
global first arrival gives an HJB value field. These are alternative physical
questions, not rival notations for one universal trajectory.

Curvature decides whether nonlinear path bending can be removed by changing
coordinates. Conformal maps preserve angles, not travel time. The earlier
complex square map makes torque motion visible through
\(\dot\zeta=2w\dot w\) but is singular at zero and does not generically flatten
the dynamic metric. Hyperbolic coordinates are more useful dynamically:
\(\dot i_T\) changes torque while \(\dot\alpha\) reallocates current at
constant torque. The metric natural gradient and tangent projector make this
task/null-space split directly implementable.

The fast--slow experiment also prevents a tempting overclaim. A field time
constant 4.41 times the stator time constant did not make naive stator-first
motion optimal. Within the tested monotonic-torque waypoint family, the
coordinated time-oriented path was 15.8\% faster than straight interpolation,
while the energy-oriented path used 66.5\% less copper energy than the time
path at a 9.7\% time penalty. Time-scale separation supplies a moving optimum
manifold and useful controller coordinates; voltage coupling still decides
the optimum.

The representation study adds a fourth, deliberately layered answer. The
rotor \(dq\) frame remains an excellent spatial interface because it freezes
sinusoidal variables and saliency. Stator-flux \(ft\) coordinates make torque
and the high-speed voltage limit direct, while active flux absorbs saliency
and excitation into one torque-producing flux factor. Neither eliminates its
observer or magnetic-model burden. At regular points, torque can itself label
a local coordinate; hyperbolic \((i_T,\alpha,z_0)\) variables then separate
torque amplitude from two EESM allocation motions. Performance coordinates
\((M,P_{\mathrm{Cu}},U^2)\) straighten three level-set families but are only
local and can be badly conditioned.

Loss-metric natural gradients and state-dependent frames expose useful local
torque and voltage directions under saturation, but their moving bases add
connection terms and need not integrate to global coordinates. After loss
whitening, the sampled torque and voltage forms do not commute, so no constant
orthogonal frame diagonalizes every objective. Approximate joint
diagonalization provides a cheap fixed compromise; feedback linearization
provides direct task-rate inputs only where its decoupling matrix is
nonsingular and converter constraints remain satisfied. The numerical study
verifies every forward/inverse identity, maps singular conditioning regions,
and shows that the conventional \(dq\) stack is still the lowest-cost baseline
in the explicit operation-count model.

The inverse-design view now includes dynamic and representation capability. Local metric
eigenvalues, finite-horizon Gramian spectra, excitation-to-stator time-scale
ratio, and minimum time or energy for representative transitions complement
steady torque--speed capability. Their intersection reveals machines that are
efficient but dynamically slow, fast but statically inadequate, or efficient
yet awkward to control. Thus the complete framework asks four linked
questions: which variables expose the physics, where should the machine
operate, how should it move between operating points, and which physical
machine makes these points and paths possible. The common philosophy remains
to seek the representation in which the physical constraint becomes the
simplest geometric object---and to use conditioning, connection terms,
curvature, residuals, and exact converter checks when no global simplification
exists.
